\documentclass[10pt,a4paper]{amsart}
\usepackage{amsmath,amssymb,amsthm,graphicx}
% Retained preamble for the Stacks Project chapter snapshot in this directory.
% The chapter's own preamble.tex is a copy of the shared Stacks Project preamble; the block
% below is sliced MECHANICALLY, line for line, from it. Nothing was added, removed or
% redefined; the wrapper class, the xr-hyper/hyperref wiring and the externaldocument list
% are omitted because the assembler supplies the document class.
%
% lines taken from preamble.tex: 3, 5, 9, 11, 12, 156, 157, 158, ... (51 lines total)
%
% package coverage restored from the chapter's own preamble (2026-10-13): xy, xy-2cell

\usepackage[all]{xy}
\xyoption{2cell}
\UseAllTwocells

\newenvironment{reference}{\comment}{\endcomment}
\newenvironment{slogan}{\comment}{\endcomment}
\newenvironment{history}{\comment}{\endcomment}
\theoremstyle{plain}
\newtheorem{theorem}[subsection]{Theorem}
\newtheorem{proposition}[subsection]{Proposition}
\newtheorem{lemma}[subsection]{Lemma}
\theoremstyle{definition}
\newtheorem{definition}[subsection]{Definition}
\newtheorem{example}[subsection]{Example}
\newtheorem{exercise}[subsection]{Exercise}
\newtheorem{situation}[subsection]{Situation}
\theoremstyle{remark}
\newtheorem{remark}[subsection]{Remark}
\newtheorem{remarks}[subsection]{Remarks}
\numberwithin{equation}{subsection}
\def\lim{\mathop{\mathrm{lim}}\nolimits}
\def\colim{\mathop{\mathrm{colim}}\nolimits}
\def\Spec{\mathop{\mathrm{Spec}}}
\def\Hom{\mathop{\mathrm{Hom}}\nolimits}
\def\Ext{\mathop{\mathrm{Ext}}\nolimits}
\def\SheafHom{\mathop{\mathcal{H}\!\mathit{om}}\nolimits}
\def\SheafExt{\mathop{\mathcal{E}\!\mathit{xt}}\nolimits}
\def\Sch{\mathit{Sch}}
\def\Mor{\mathop{\mathrm{Mor}}\nolimits}
\def\Ob{\mathop{\mathrm{Ob}}\nolimits}
\def\Sh{\mathop{\mathit{Sh}}\nolimits}
\def\NL{\mathop{N\!L}\nolimits}
\def\CH{\mathop{\mathrm{CH}}\nolimits}
\def\proetale{{pro\text{-}\acute{e}tale}}
\def\etale{{\acute{e}tale}}
\def\QCoh{\mathit{QCoh}}
\def\Ker{\mathop{\mathrm{Ker}}}
\def\Im{\mathop{\mathrm{Im}}}
\def\Coker{\mathop{\mathrm{Coker}}}
\def\Coim{\mathop{\mathrm{Coim}}}
\DeclareMathSymbol{\boxtimes}{\mathbin}{AMSa}{"02}
\def\QCohstack{\mathcal{QC}\!\mathit{oh}}
\def\Cohstack{\mathcal{C}\!\mathit{oh}}
\def\Spacesstack{\mathcal{S}\!\mathit{paces}}
\def\Quotfunctor{\mathrm{Quot}}
\def\Hilbfunctor{\mathrm{Hilb}}
\def\Curvesstack{\mathcal{C}\!\mathit{urves}}
\def\Polarizedstack{\mathcal{P}\!\mathit{olarized}}
\def\Complexesstack{\mathcal{C}\!\mathit{omplexes}}
\def\Pic{\mathop{\mathrm{Pic}}\nolimits}
\def\Picardstack{\mathcal{P}\!\mathit{ic}}
\def\Picardfunctor{\mathrm{Pic}}
\def\Deformationcategory{\mathcal{D}\!\mathit{ef}}
\usepackage{cleveref}
\crefname{originallemma}{original lemma}{original lemmas}
\Crefname{originallemma}{Original lemma}{Original lemmas}
\makeatletter
\newlabel{categories-lemma-equivalence-fibred-categories}{{003Z}{source}{}{sourcepointer-1}{}}
\newlabel{categories-lemma-equivalence-fibred-categories@cref}{{[originallemma][0][]003Z}{[1][0][]source}{}{}{}}
\makeatother
\crefname{originallemma}{original lemma}{original lemmas}
\Crefname{originallemma}{Original lemma}{Original lemmas}
\makeatletter
\newlabel{definition-RS}{{07WP}{source}{}{sourcepointer-2}{}}
\newlabel{definition-RS@cref}{{[originallemma][0][]07WP}{[1][0][]source}{}{}{}}
\makeatother
\crefname{originallemma}{original lemma}{original lemmas}
\Crefname{originallemma}{Original lemma}{Original lemmas}
\makeatletter
\newlabel{equation-factor}{{07T3}{source}{}{sourcepointer-3}{}}
\newlabel{equation-factor@cref}{{[originallemma][0][]07T3}{[1][0][]source}{}{}{}}
\makeatother
\crefname{originallemma}{original lemma}{original lemmas}
\Crefname{originallemma}{Original lemma}{Original lemmas}
\makeatletter
\newlabel{formal-defos-situation-change-of-fields}{{07W8}{source}{}{sourcepointer-4}{}}
\newlabel{formal-defos-situation-change-of-fields@cref}{{[originallemma][0][]07W8}{[1][0][]source}{}{}{}}
\makeatother
\crefname{originallemma}{original lemma}{original lemmas}
\Crefname{originallemma}{Original lemma}{Original lemmas}
\makeatletter
\newlabel{more-morphisms-lemma-pushout-along-thickening}{{07RT}{source}{}{sourcepointer-5}{}}
\newlabel{more-morphisms-lemma-pushout-along-thickening@cref}{{[originallemma][0][]07RT}{[1][0][]source}{}{}{}}
\makeatother
\title{Field change identifies deformation categories under (RS)}
\author{The Stacks Project Authors}
\date{stacks-project:artin.tex:local-snapshot-20261012; GNU FDL 1.2 or later}
\begin{document}
\maketitle
\noindent\textit{Editorial note.} Statement span is lines 902-917 of sources/stacks-artin/pinned-original-source.tex; proof span is lines 919-972, the complete author proof (52 body lines).\par\medskip
\noindent\textit{Editorial note.} Navigation stubs: 5 cited-but-not-included labels are supplied as non-load-bearing pointers with their pinned tags where the snapshot tag index has them.\par\medskip
\noindent\textit{Editorial note.} extra\_packages is empty on purpose: the retained preamble already loads xy with options and a second declaration clashes in the assembler wrapper.\par\medskip
\noindent\textit{Editorial note.} Text is retained verbatim, including the author's own typos and placeholders.\par\medskip
\noindent\textit{Editorial note.} Queue index 24 of queue\_artin.json (pb=1766, stmt 902-917, proof 919-972).\par\medskip
\noindent\textit{Editorial note.} stacks-artin wave ar01 (2026-10-14): fresh adjudication accept per lane-s3/artin\_judgments.md; pool per Chief 10:5x direct dispatch (idle-first over S12's 01:5x note, SMC precedent); artin.tex 99 rows/67 with proof judged 46 accept / 21 reject / 5 AGP-overlap avoidance rows; patch-sd-envs contract spans scanned.\par\medskip
\noindent\textit{Allowed standard prerequisites (editorial).}\begin{itemize}
\item Explicit prerequisite boundary. The cited inputs are prerequisites; the argument in the retained span is not replaced by any cited result.
\end{itemize}\medskip
\ifdefined\hypertarget\hypertarget{sourcepointer-1}{}\fi
\noindent\textit{Source reference.} Original lemma 003Z: sources/stacks-artin/pinned-original-source.tex, source lines 972--972. Cited by the retained statement or proof but not part of the collected spans; supplied only as a navigation pointer. The argument does not depend on its content.\par\medskip
\ifdefined\hypertarget\hypertarget{sourcepointer-2}{}\fi
\noindent\textit{Source reference.} Original lemma 07WP: sources/stacks-artin/pinned-original-source.tex, source lines 710--725. Cited by the retained statement or proof but not part of the collected spans; supplied only as a navigation pointer. The argument does not depend on its content.\par\medskip
\ifdefined\hypertarget\hypertarget{sourcepointer-3}{}\fi
\noindent\textit{Source reference.} Original lemma 07T3: sources/stacks-artin/pinned-original-source.tex, source lines 441--443. Cited by the retained statement or proof but not part of the collected spans; supplied only as a navigation pointer. The argument does not depend on its content.\par\medskip
\ifdefined\hypertarget\hypertarget{sourcepointer-4}{}\fi
\noindent\textit{Source reference.} Original lemma 07W8: sources/stacks-artin/pinned-original-source.tex, source lines 972--972. Cited by the retained statement or proof but not part of the collected spans; supplied only as a navigation pointer. The argument does not depend on its content.\par\medskip
\ifdefined\hypertarget\hypertarget{sourcepointer-5}{}\fi
\noindent\textit{Source reference.} Original lemma 07RT: sources/stacks-artin/pinned-original-source.tex, source lines 972--972. Cited by the retained statement or proof but not part of the collected spans; supplied only as a navigation pointer. The argument does not depend on its content.\par\medskip
\begin{lemma}
\label{lemma-change-of-field}
Let $S$ be a locally Noetherian scheme. Let $\mathcal{X}$ be a category
fibred in groupoids over $(\Sch/S)_{fppf}$. Let $k$ be a
field of finite type over $S$ and let $l/k$ be a finite extension.
Let $x_0$ be an object of $\mathcal{F}$ lying over $\Spec(k)$.
Denote $x_{l, 0}$ the restriction of $x_0$ to $\Spec(l)$.
Then there is a canonical functor
$$
(\mathcal{F}_{\mathcal{X}, k , x_0})_{l/k}
\longrightarrow
\mathcal{F}_{\mathcal{X}, l, x_{l, 0}}
$$
of categories cofibred in groupoids over $\mathcal{C}_{\Lambda, l}$.
If $\mathcal{X}$ satisfies (RS), then this functor is an equivalence.
\end{lemma}

\begin{proof}
Consider a factorization
$$
\Spec(l) \to \Spec(B) \to S
$$
as in (\ref{equation-factor}). By definition we have
$$
(\mathcal{F}_{\mathcal{X}, k , x_0})_{l/k}(B) =
\mathcal{F}_{\mathcal{X}, k, x_0}(B \times_l k)
$$
see Formal Deformation Theory, Situation
\ref{formal-defos-situation-change-of-fields}. Thus an object of this
is a morphism $x_0 \to x$ of $\mathcal{X}$ lying over the morphism
$\Spec(k) \to \Spec(B \times_l k)$. Choosing pullback functor for $\mathcal{X}$
we can associate to $x_0 \to x$ the morphism $x_{l, 0} \to x_B$
where $x_B$ is the restriction of $x$ to $\Spec(B)$ (via the morphism
$\Spec(B) \to \Spec(B \times_l k)$ coming from $B \times_l k \subset B$).
This construction is functorial in $B$ and compatible with morphisms.

\medskip\noindent
Next, assume $\mathcal{X}$ satisfies (RS). Consider the diagrams
$$
\vcenter{
\xymatrix{
l & B \ar[l] \\
k \ar[u] & B \times_l k \ar[l] \ar[u]
}
}
\quad\text{and}\quad
\vcenter{
\xymatrix{
\Spec(l) \ar[d] \ar[r] & \Spec(B) \ar[d] \\
\Spec(k) \ar[r] & \Spec(B \times_l k)
}
}
$$
The diagram on the left is a fibre product of rings. The diagram on the
right is a pushout in the category of schemes, see
More on Morphisms, Lemma \ref{more-morphisms-lemma-pushout-along-thickening}.
These schemes are all of finite type over $S$ (see remarks following
Definition \ref{definition-RS}). Hence (RS) kicks in to give an equivalence
of fibre categories
$$
\mathcal{X}_{\Spec(B \times_l k)}
\longrightarrow
\mathcal{X}_{\Spec(k)}
\times_{\mathcal{X}_{\Spec(l)}}
\mathcal{X}_{\Spec(B)}
$$
This implies that the functor defined above gives an equivalence of
fibre categories. Hence the functor is an equivalence on categories
cofibred in groupoids by (the dual of)
Categories, Lemma \ref{categories-lemma-equivalence-fibred-categories}.
\end{proof}

\end{document}
